
\chapter{Giới thiệu}

\section{Mục đích}

Tài liệu Đặc tả yêu cầu phần mềm (Software Requirements Specification - SRS) được xây dựng nhằm mô tả các yêu cầu chức năng và phi chức năng của hệ thống \textbf{Quản lý dịch vụ bảo trì máy lọc nước}. Tài liệu đóng vai trò là cơ sở cho các giai đoạn phân tích, thiết kế, phát triển, kiểm thử và bảo trì hệ thống.

Ngoài ra, SRS còn là tài liệu tham chiếu giúp các bên liên quan có cùng cách hiểu về phạm vi dự án, chức năng của hệ thống cũng như các ràng buộc kỹ thuật trước khi tiến hành hiện thực.

Do dự án được thực hiện với mục đích thực tập và xây dựng hệ thống minh họa (Demo System), các yêu cầu được tổng hợp dựa trên quy trình nghiệp vụ phổ biến của dịch vụ bảo trì máy lọc nước thay vì xuất phát từ tài liệu đặc tả của một doanh nghiệp cụ thể.

\section{Phạm vi}

Hệ thống Quản lý dịch vụ bảo trì máy lọc nước là một ứng dụng Web hỗ trợ quản lý thông tin sản phẩm, khách hàng, thiết bị, hợp đồng bảo trì, tồn kho và hoạt động của doanh nghiệp trên một nền tảng tập trung.

Các chức năng chính của hệ thống bao gồm:

\begin{itemize}
    \item Dashboard thống kê tổng quan.
    \item Quản lý sản phẩm.
    \item Quản lý tồn kho (Inventory).
    \item Quản lý lô hàng (Batch).
    \item Quản lý khách hàng.
    \item Quản lý hợp đồng.
    \item Quản lý thiết bị.
    \item Dashboard theo dõi dữ liệu thiết bị.
    \item Quản lý nhân viên.
    \item Quản lý Activity Log.
    \item Quản lý thông tin cá nhân.
    \item Đăng nhập và phân quyền người dùng.
\end{itemize}


Các đối tượng trên sử dụng tài liệu này để hiểu yêu cầu của hệ thống, thống nhất phạm vi phát triển và làm cơ sở cho việc thiết kế, hiện thực cũng như kiểm thử phần mềm.

\section{Thuật ngữ và từ viết tắt}

\begin{table}[H]
\centering
\caption{Thuật ngữ và từ viết tắt}
\begin{tabularx}{\textwidth}{|l|X|}
\hline
\textbf{Thuật ngữ} & \textbf{Mô tả} \\
\hline
SRS & Software Requirements Specification - Tài liệu đặc tả yêu cầu phần mềm. \\
\hline
UI & User Interface - Giao diện người dùng. \\
\hline
UX & User Experience - Trải nghiệm người dùng. \\
\hline
CRUD & Create, Read, Update, Delete. \\
\hline
Dashboard & Giao diện hiển thị dữ liệu tổng hợp bằng biểu đồ hoặc thống kê. \\
\hline
Inventory & Danh sách thiết bị đang lưu trong kho. \\
\hline
Batch & Lô hàng nhập kho. \\
\hline
Device & Thiết bị máy lọc nước đã được quản lý trong hệ thống. \\
\hline
Activity Log & Nhật ký ghi nhận các thao tác của người dùng. \\
\hline
IoT & Internet of Things. \\
\hline
MQTT & Giao thức truyền dữ liệu giữa thiết bị IoT và Server. \\
\hline
API & Application Programming Interface. \\
\hline
\end{tabularx}
\end{table}

\section{Cấu trúc tài liệu}

Tài liệu được tổ chức thành các chương như sau:

\begin{itemize}
    \item \textbf{Chương 1}: Giới thiệu mục đích, phạm vi, thuật ngữ và cấu trúc của tài liệu.
    
    \item \textbf{Chương 2}: Mô tả tổng quan hệ thống, các nhóm người dùng và môi trường hoạt động.
    
    \item \textbf{Chương 3}: Mô tả các giao diện bên ngoài của hệ thống bao gồm giao diện người dùng, phần cứng, phần mềm và giao tiếp.
    
    \item \textbf{Chương 4}: Đặc tả các yêu cầu chức năng của từng module trong hệ thống.
    
    \item \textbf{Chương 5}: Đặc tả các yêu cầu phi chức năng như hiệu năng, bảo mật, khả năng mở rộng và tính ổn định.
    
    \item \textbf{Chương 6}: Đặc tả các Use Case chính của hệ thống.
    \item \textbf{Chương 7}: mô tả các yêu cầu về dữ liệu của hệ thống.
\end{itemize}

